%!TeX root=../emmatop.tex
\chapter[Chapter \thechapter]{}
\lettrine[lraise=0.3]{I}{t} may be possible to do without dancing entirely. Instances have been known of young people passing many, many months successively, without being at any ball of any description, and no material injury accrue either to body or mind;—but when a beginning is made—when the felicities of rapid motion have once been, though slightly, felt—it must be a very heavy set that does not ask for more.

Frank Churchill had danced once at Highbury, and longed to dance again; and the last half-hour of an evening which Mr Woodhouse was persuaded to spend with his daughter at Randalls, was passed by the two young people in schemes on the subject. Frank's was the first idea; and his the greatest zeal in pursuing it; for the lady was the best judge of the difficulties, and the most solicitous for accommodation and appearance. But still she had inclination enough for shewing people again how delightfully Mr Frank Churchill and Miss Woodhouse danced—for doing that in which she need not blush to compare herself with Jane Fairfax—and even for simple dancing itself, without any of the wicked aids of vanity—to assist him first in pacing out the room they were in to see what it could be made to hold—and then in taking the dimensions of the other parlour, in the hope of discovering, in spite of all that Mr Weston could say of their exactly equal size, that it was a little the largest.

His first proposition and request, that the dance begun at Mr Cole's should be finished there—that the same party should be collected, and the same musician engaged, met with the readiest acquiescence. Mr Weston entered into the idea with thorough enjoyment, and Mrs Weston most willingly undertook to play as long as they could wish to dance; and the interesting employment had followed, of reckoning up exactly who there would be, and portioning out the indispensable division of space to every couple.

<You and Miss Smith, and Miss Fairfax, will be three, and the two Miss Coxes five,> had been repeated many times over. <And there will be the two Gilberts, young Cox, my father, and myself, besides Mr Knightley. Yes, that will be quite enough for pleasure. You and Miss Smith, and Miss Fairfax, will be three, and the two Miss Coxes five; and for five couple there will be plenty of room.>

But soon it came to be on one side,

<But will there be good room for five couple?—I really do not think there will.>

On another,

<And after all, five couple are not enough to make it worth while to stand up. Five couple are nothing, when one thinks seriously about it. It will not do to invite five couple. It can be allowable only as the thought of the moment.>

Somebody said that Miss Gilbert was expected at her brother's, and must be invited with the rest. Somebody else believed Mrs Gilbert would have danced the other evening, if she had been asked. A word was put in for a second young Cox; and at last, Mr Weston naming one family of cousins who must be included, and another of very old acquaintance who could not be left out, it became a certainty that the five couple would be at least ten, and a very interesting speculation in what possible manner they could be disposed of.

The doors of the two rooms were just opposite each other. <Might not they use both rooms, and dance across the passage?> It seemed the best scheme; and yet it was not so good but that many of them wanted a better. Emma said it would be awkward; Mrs Weston was in distress about the supper; and Mr Woodhouse opposed it earnestly, on the score of health. It made him so very unhappy, indeed, that it could not be persevered in.

<Oh! no,> said he; <it would be the extreme of imprudence. I could not bear it for Emma!—Emma is not strong. She would catch a dreadful cold. So would poor little Harriet. So you would all. Mrs Weston, you would be quite laid up; do not let them talk of such a wild thing. Pray do not let them talk of it. That young man (speaking lower) is very thoughtless. Do not tell his father, but that young man is not quite the thing. He has been opening the doors very often this evening, and keeping them open very inconsiderately. He does not think of the draught. I do not mean to set you against him, but indeed he is not quite the thing!>

Mrs Weston was sorry for such a charge. She knew the importance of it, and said every thing in her power to do it away. Every door was now closed, the passage plan given up, and the first scheme of dancing only in the room they were in resorted to again; and with such good-will on Frank Churchill's part, that the space which a quarter of an hour before had been deemed barely sufficient for five couple, was now endeavoured to be made out quite enough for ten.

<We were too magnificent,> said he. <We allowed unnecessary room. Ten couple may stand here very well.>

Emma demurred. <It would be a crowd—a sad crowd; and what could be worse than dancing without space to turn in?>

<Very true,> he gravely replied; <it was very bad.> But still he went on measuring, and still he ended with,

<I think there will be very tolerable room for ten couple.>

<No, no,> said she, <you are quite unreasonable. It would be dreadful to be standing so close! Nothing can be farther from pleasure than to be dancing in a crowd—and a crowd in a little room!>

<There is no denying it,> he replied. <I agree with you exactly. A crowd in a little room—Miss Woodhouse, you have the art of giving pictures in a few words. Exquisite, quite exquisite!—Still, however, having proceeded so far, one is unwilling to give the matter up. It would be a disappointment to my father—and altogether—I do not know that—I am rather of opinion that ten couple might stand here very well.>

Emma perceived that the nature of his gallantry was a little self-willed, and that he would rather oppose than lose the pleasure of dancing with her; but she took the compliment, and forgave the rest. Had she intended ever to marry him, it might have been worth while to pause and consider, and try to understand the value of his preference, and the character of his temper; but for all the purposes of their acquaintance, he was quite amiable enough.

Before the middle of the next day, he was at Hartfield; and he entered the room with such an agreeable smile as certified the continuance of the scheme. It soon appeared that he came to announce an improvement.

<Well, Miss Woodhouse,> he almost immediately began, <your inclination for dancing has not been quite frightened away, I hope, by the terrors of my father's little rooms. I bring a new proposal on the subject:—a thought of my father's, which waits only your approbation to be acted upon. May I hope for the honour of your hand for the two first dances of this little projected ball, to be given, not at Randalls, but at the Crown Inn?>

<The Crown!>

<Yes; if you and Mr Woodhouse see no objection, and I trust you cannot, my father hopes his friends will be so kind as to visit him there. Better accommodations, he can promise them, and not a less grateful welcome than at Randalls. It is his own idea. Mrs Weston sees no objection to it, provided you are satisfied. This is what we all feel. Oh! you were perfectly right! Ten couple, in either of the Randalls rooms, would have been insufferable!—Dreadful!—I felt how right you were the whole time, but was too anxious for securing any thing to like to yield. Is not it a good exchange?—You consent—I hope you consent?>

<It appears to me a plan that nobody can object to, if Mr and Mrs Weston do not. I think it admirable; and, as far as I can answer for myself, shall be most happy—It seems the only improvement that could be. Papa, do you not think it an excellent improvement?>

She was obliged to repeat and explain it, before it was fully comprehended; and then, being quite new, farther representations were necessary to make it acceptable.

<No; he thought it very far from an improvement—a very bad plan—much worse than the other. A room at an inn was always damp and dangerous; never properly aired, or fit to be inhabited. If they must dance, they had better dance at Randalls. He had never been in the room at the Crown in his life—did not know the people who kept it by sight.—Oh! no—a very bad plan. They would catch worse colds at the Crown than anywhere.>

<I was going to observe, sir,> said Frank Churchill, <that one of the great recommendations of this change would be the very little danger of any body's catching cold—so much less danger at the Crown than at Randalls! Mr Perry might have reason to regret the alteration, but nobody else could.>

<Sir,> said Mr Woodhouse, rather warmly, <you are very much mistaken if you suppose Mr Perry to be that sort of character. Mr Perry is extremely concerned when any of us are ill. But I do not understand how the room at the Crown can be safer for you than your father's house.>

<From the very circumstance of its being larger, sir. We shall have no occasion to open the windows at all—not once the whole evening; and it is that dreadful habit of opening the windows, letting in cold air upon heated bodies, which (as you well know, sir) does the mischief.>

<Open the windows!—but surely, Mr Churchill, nobody would think of opening the windows at Randalls. Nobody could be so imprudent! I never heard of such a thing. Dancing with open windows!—I am sure, neither your father nor Mrs Weston (poor Miss Taylor that was) would suffer it.>

<Ah! sir—but a thoughtless young person will sometimes step behind a window-curtain, and throw up a sash, without its being suspected. I have often known it done myself.>

<Have you indeed, sir?—Bless me! I never could have supposed it. But I live out of the world, and am often astonished at what I hear. However, this does make a difference; and, perhaps, when we come to talk it over—but these sort of things require a good deal of consideration. One cannot resolve upon them in a hurry. If Mr and Mrs Weston will be so obliging as to call here one morning, we may talk it over, and see what can be done.>

<But, unfortunately, sir, my time is so limited\longdash>

<Oh!> interrupted Emma, <there will be plenty of time for talking every thing over. There is no hurry at all. If it can be contrived to be at the Crown, papa, it will be very convenient for the horses. They will be so near their own stable.>

<So they will, my dear. That is a great thing. Not that James ever complains; but it is right to spare our horses when we can. If I could be sure of the rooms being thoroughly aired—but is Mrs Stokes to be trusted? I doubt it. I do not know her, even by sight.>

<I can answer for every thing of that nature, sir, because it will be under Mrs Weston's care. Mrs Weston undertakes to direct the whole.>

<There, papa!—Now you must be satisfied—Our own dear Mrs Weston, who is carefulness itself. Do not you remember what Mr Perry said, so many years ago, when I had the measles? <If Miss Taylor undertakes to wrap Miss Emma up, you need not have any fears, sir.> How often have I heard you speak of it as such a compliment to her!>

<Aye, very true. Mr Perry did say so. I shall never forget it. Poor little Emma! You were very bad with the measles; that is, you would have been very bad, but for Perry's great attention. He came four times a day for a week. He said, from the first, it was a very good sort—which was our great comfort; but the measles are a dreadful complaint. I hope whenever poor Isabella's little ones have the measles, she will send for Perry.>

<My father and Mrs Weston are at the Crown at this moment,> said Frank Churchill, <examining the capabilities of the house. I left them there and came on to Hartfield, impatient for your opinion, and hoping you might be persuaded to join them and give your advice on the spot. I was desired to say so from both. It would be the greatest pleasure to them, if you could allow me to attend you there. They can do nothing satisfactorily without you.>

Emma was most happy to be called to such a council; and her father, engaging to think it all over while she was gone, the two young people set off together without delay for the Crown. There were Mr and Mrs Weston; delighted to see her and receive her approbation, very busy and very happy in their different way; she, in some little distress; and he, finding every thing perfect.

<Emma,> said she, <this paper is worse than I expected. Look! in places you see it is dreadfully dirty; and the wainscot is more yellow and forlorn than any thing I could have imagined.>

<My dear, you are too particular,> said her husband. <What does all that signify? You will see nothing of it by candlelight. It will be as clean as Randalls by candlelight. We never see any thing of it on our club-nights.>

The ladies here probably exchanged looks which meant, <Men never know when things are dirty or not;> and the gentlemen perhaps thought each to himself, <Women will have their little nonsenses and needless cares.>

One perplexity, however, arose, which the gentlemen did not disdain. It regarded a supper-room. At the time of the ballroom's being built, suppers had not been in question; and a small card-room adjoining, was the only addition. What was to be done? This card-room would be wanted as a card-room now; or, if cards were conveniently voted unnecessary by their four selves, still was it not too small for any comfortable supper? Another room of much better size might be secured for the purpose; but it was at the other end of the house, and a long awkward passage must be gone through to get at it. This made a difficulty. Mrs Weston was afraid of draughts for the young people in that passage; and neither Emma nor the gentlemen could tolerate the prospect of being miserably crowded at supper.

Mrs Weston proposed having no regular supper; merely sandwiches, \&c., set out in the little room; but that was scouted as a wretched suggestion. A private dance, without sitting down to supper, was pronounced an infamous fraud upon the rights of men and women; and Mrs Weston must not speak of it again. She then took another line of expediency, and looking into the doubtful room, observed,

<I do not think it is so very small. We shall not be many, you know.>

And Mr Weston at the same time, walking briskly with long steps through the passage, was calling out,

<You talk a great deal of the length of this passage, my dear. It is a mere nothing after all; and not the least draught from the stairs.>

<I wish,> said Mrs Weston, <one could know which arrangement our guests in general would like best. To do what would be most generally pleasing must be our object—if one could but tell what that would be.>

<Yes, very true,> cried Frank, <very true. You want your neighbours' opinions. I do not wonder at you. If one could ascertain what the chief of them—the Coles, for instance. They are not far off. Shall I call upon them? Or Miss Bates? She is still nearer.—And I do not know whether Miss Bates is not as likely to understand the inclinations of the rest of the people as any body. I think we do want a larger council. Suppose I go and invite Miss Bates to join us?>

<Well—if you please,> said Mrs Weston rather hesitating, <if you think she will be of any use.>

<You will get nothing to the purpose from Miss Bates,> said Emma. <She will be all delight and gratitude, but she will tell you nothing. She will not even listen to your questions. I see no advantage in consulting Miss Bates.>

<But she is so amusing, so extremely amusing! I am very fond of hearing Miss Bates talk. And I need not bring the whole family, you know.>

Here Mr Weston joined them, and on hearing what was proposed, gave it his decided approbation.

<Aye, do, Frank.—Go and fetch Miss Bates, and let us end the matter at once. She will enjoy the scheme, I am sure; and I do not know a properer person for shewing us how to do away difficulties. Fetch Miss Bates. We are growing a little too nice. She is a standing lesson of how to be happy. But fetch them both. Invite them both.>

<Both sir! Can the old lady?>...

<The old lady! No, the young lady, to be sure. I shall think you a great blockhead, Frank, if you bring the aunt without the niece.>

<Oh! I beg your pardon, sir. I did not immediately recollect. Undoubtedly if you wish it, I will endeavour to persuade them both.> And away he ran.

Long before he reappeared, attending the short, neat, brisk-moving aunt, and her elegant niece,—Mrs Weston, like a sweet-tempered woman and a good wife, had examined the passage again, and found the evils of it much less than she had supposed before—indeed very trifling; and here ended the difficulties of decision. All the rest, in speculation at least, was perfectly smooth. All the minor arrangements of table and chair, lights and music, tea and supper, made themselves; or were left as mere trifles to be settled at any time between Mrs Weston and Mrs Stokes.—Every body invited, was certainly to come; Frank had already written to Enscombe to propose staying a few days beyond his fortnight, which could not possibly be refused. And a delightful dance it was to be.

\begin{figure}[tbph]
\centering
\includegraphics[width=.8\linewidth]{29asked}
\caption{<He has asked her, my dear.>}
\end{figure}

Most cordially, when Miss Bates arrived, did she agree that it must. As a counsellor she was not wanted; but as an approver, (a much safer character,) she was truly welcome. Her approbation, at once general and minute, warm and incessant, could not but please; and for another half-hour they were all walking to and fro, between the different rooms, some suggesting, some attending, and all in happy enjoyment of the future. The party did not break up without Emma's being positively secured for the two first dances by the hero of the evening, nor without her overhearing Mr Weston whisper to his wife, <He has asked her, my dear. That's right. I knew he would!>